%%%%%%%%%%%% STRUCTURE %%%%%%%%%%%%%%%
\documentclass[a4paper,12pt]{article}
\usepackage[T1]{fontenc}
\usepackage[utf8]{inputenc}
\usepackage[brazil]{babel}
\usepackage{lmodern}
\usepackage{setspace}
\usepackage[top=2cm, bottom=2cm, left=2cm, right=2cm]{geometry}
\usepackage{listings}
\usepackage{xcolor}
\usepackage{amsmath}
\usepackage{ragged2e}
\usepackage{float}

\lstset{
    basicstyle=\ttfamily\small,
    keywordstyle=\color{blue}\bfseries,
    stringstyle=\color{red},
    commentstyle=\color{gray}\itshape,
    numbers=left,
    numberstyle=\tiny\color{gray},
    stepnumber=1,
    numbersep=8pt,
    backgroundcolor=\color{white},
    showspaces=false,
    showstringspaces=false,
    showtabs=false,
    frame=single,
    rulecolor=\color{black!20},
    tabsize=4,
    captionpos=b,
    breaklines=true,
    breakatwhitespace=false,
}
%%%%%%%%%%%%%%%%%%%%%%%%%%%%%%%%%%%%%%

%%%%%%%%%%%%%%%% PAGES STYLE %%%%%%%%%
\usepackage{fancyhdr}
\setlength{\headheight}{15pt}
\fancypagestyle{main}{
\renewcommand{\headrulewidth}{0pt}
\fancyhead[RO]{\thepage}
\fancyfoot[CO]{}
}
%%%%%%%%%%%%%%%%%%%%%%%%%%%%%%%%%%%%%%

\usepackage{graphicx}
\usepackage{epstopdf}
\usepackage{subfig}
\usepackage{mathptmx}
\usepackage{changepage}
%\usepackage[alf]{abntex2cite}

%%%%%%%%%%% PDF METADATA %%%%%%%%%%%%%
\usepackage[ pdftitle={RELATÓRIO SISTEMAS DE CONTROLE},
pdfsubject={LABORATÓRIO DE CONTROLE},
pdfkeywords={Controle,Automação,UFRN,DCA},
hidelinks]{hyperref}
%%%%%%%%%%%%%%%%%%%%%%%%%%%%%%%%%%%%%%

\begin{document}

\onehalfspacing

\thispagestyle{empty}

\setcounter{page}{1}

%%%%%%%%%%%% LOGOS %%%%%%%%%%%%%%%%%%%

\begin{figure}[!ht]

\centering

\subfloat{
\includegraphics[width=2.7cm]{img/UFRN-eps-converted-to.pdf}
\label{UFRN Logo}
}
\hspace{11.09cm}
\subfloat{
\includegraphics[width=2.4cm]{img/DCA-eps-converted-to.pdf}
\label{DCA Logo}
}

%\caption{}
\label{Logos}

\end{figure}

%%%%%%%%%%%%%%% CAPA %%%%%%%%%%%%%%%%%

\vspace{-1cm}

\begin{center}
{\bf{\normalsize UNIVERSIDADE FEDERAL DO RIO GRANDE DO NORTE\\
CENTRO DE TECNOLOGIA\\
DEPARTAMENTO DE ENGENHARIA DE COMPUTAÇÃO E AUTOMAÇÃO\\
ENGENHARIA DE COMPUTAÇÃO E ENGENHARIA MECATRÔNICA
}}


\vspace{3.6cm}

{\bf{\large RELATÓRIO DA 2º UNIDADE\\
PROJETO DE CONTROLADORES PROPORCIONAL, PROPORCIONAL-DERIVATIVO, PROPORCIONAL-INTEGRATIVO E PROPORCIONAL-INTEGRATIVO-DERIVATIVO\\
}}
\vspace{1.5cm}
{\large TURMA: 1\\
	GRUPO 20262G4T1}

\vspace{3.6cm}



\begin{flushright}
\begin{normalsize}
Danilo Isamu Pereira Hikague: 20230072008\\
\vspace{0.8cm}
Murilo de Lima Barros: 20250070735\\
\vspace{0.8cm}
Rafael Cabral Melo: 20210080441\\
\vspace{0.8cm}
Ramon Vinicius Ferreira de Souza: 20220017441\\
\vspace{0.8cm}
Virna Emanuelle Veloso Aguiar:: 20220025354\\
\end{normalsize}
\end{flushright}


\vspace{2.5cm}

{\large Natal-RN\\
2026}

\end{center}

\newpage

%%%%%%%%%%%%%%%  CONTRA-CAPA %%%%%%%%%

\thispagestyle{empty}

\begin{center}
\begin{normalsize}
Danilo Isamu Pereira Hikague: 20230072008\\
\vspace{0.8cm}
Murilo de Lima Barros: 20250070735\\
\vspace{0.8cm}
Rafael Cabral Melo: 20210080441\\
\vspace{0.8cm}
Ramon Vinicius Ferreira de Souza: 20220017441\\
\vspace{0.8cm}
Virna Emanuelle Veloso Aguiar:: 20220025354\\

\end{normalsize}
\end{center}
\vspace{3cm}

{\bf{\large {\centering PROJETO DE CONTROLADORES PROPORCIONAL, PROPORCIONAL-DERIVATIVO, PROPORCIONAL-INTEGRATIVO E PROPORCIONAL-INTEGRATIVO-DERIVATIVO\\
}}}

\vspace{4cm}

\begin{adjustwidth}{7.5cm}{0cm}

{\normalsize

Segundo Relatório Parcial apresentado à disciplina de
Projeto de Sistemas de Controle - Laboratório, correspondente à
avaliação da 2º unidade do semestre 2026.2 do 7º período
dos cursos: Engenharia de Computação e Engenharia Mecatrônica da
Universidade Federal do Rio Grande do Norte, sob
orientação do {\bf Prof. Fábio Meneghetti Ugulino de
Araújo.}

}

\end{adjustwidth}

\vspace{2cm}

\begin{center}

Professor:  Fábio Meneghetti Ugulino de Araújo.

\vspace{2.5cm}

{\large Natal-RN\\
2026}

\end{center}

\newpage

%%%%%%%%%%%%%%%  RESUMO %%%%%%%%%%%%%%

\thispagestyle{empty}

\begin{center}
{\large \textbf{RESUMO}}
\end{center}

\vspace{3cm}

\vspace{\baselineskip}

\begin{justify}
    Resumo
\end{justify}

\vspace{\baselineskip}

\noindent\textbf{Palavras-chave:} Sistemas de controle. Tanques acoplados. Calibração de sensores. PID.

\newpage

%%%%%%%%%%%%%%%  LISTA DE SÍMBOLOS %%%

% \thispagestyle{empty}
%
% \begin{center}
% {\large \textbf{LISTA DE SÍMBOLOS}}
% \end{center}
%
% \vspace{3cm}
%
% \begin{tabular}{ l l }
% A\hspace{1.5cm} & Matriz triangular superior com diagonal unitária.\\
% \phantom{a} & \phantom{a}\\
% D\hspace{1.5cm} & Matriz diagonal obtida a partir de $W^{T}W$\\
% \phantom{a} & \phantom{a}\\
% $\theta$\hspace{1.5cm} & Vetor de parâmetros.\\
% \phantom{a} & \phantom{a}\\
% $\Xi$\hspace{1.5cm} & Vetor de resíduos de modelagem.\\
% \phantom{a} & \phantom{a}\\
% d\hspace{1.5cm} & Tempo de retardo de um sistema ou tempo morto.\\
% \phantom{a} & \phantom{a}\\
% e(k)\hspace{1.5cm} & Resíduo (Erro de Estimação mais o Ruído).\\
% \end{tabular}
%
% \newpage

%%%% LISTA DE ABREVIATURAS E SIGLAS %%

% \thispagestyle{empty}
%
% \begin{center}
% {\large \textbf{LISTA DE ABREVIATURAS E SIGLAS}}
% \end{center}
%
% \vspace{3cm}
%
% \begin{tabular}{ l l }
% ARX\hspace{1.5cm} & Matriz triangular superior com diagonal unitária.\\
% ARMAX\hspace{1.5cm} & Matriz diagonal obtida a partir de $W^{T}W$\\
% NARX\hspace{1.5cm}&Vetor de parâmetros.\\
% NARMAX\hspace{1.5cm}&Vetor de resíduos de modelagem.\\
% MQ\hspace{1.5cm}&Tempo de retardo de um sistema ou tempo morto.\\
% \end{tabular}
%
% \newpage

%%%%%%%%% LISTA DE FIGURAS %%%%%%%%%%%

\thispagestyle{empty}

\begin{center}
\listoffigures
\end{center}

\newpage

%%%%%%%%%%%%%%% SUMÁRIO %%%%%%%%%%%%%%

\thispagestyle{empty}

\begin{center}
\tableofcontents
\end{center}

\newpage

%%%%%%%%%%%%%%% INTRODUÇÃO %%%%%%%%%%%

\thispagestyle{main}

\section{INTRODUÇÃO}

\begin{justify}
    Introdução
\end{justify}

\newpage

%%%%%%% FUNDAMENTAÇÃO TEÓRICA %%%%%%%

\section{FUNDAMENTAÇÃO TEÓRICA}

\subsection{(Experimento 2A)}

\begin{justify}
    Experimento 2A
\end{justify}

\newpage

%%%%%%%%%% METODOLOGIA %%%%%%%%%%%%%%%

\thispagestyle{main}

\section{METODOLOGIA}

\begin{justify}
    Metodologia
\end{justify}

\newpage

%%%%%%%%%% RESULTADOS %%%%%%%%%%%%%%%

\thispagestyle{main}

\section{RESULTADOS}

\subsection{Experimento 2A}

\begin{justify}
    Resultados
\end{justify}

\newpage

%%%%%%%%%% CONCLUSÃO %%%%%%%%%%%%%%%

\thispagestyle{main}

\section{CONCLUSÃO}

\begin{justify}
    Conclusão
\end{justify}

\newpage

%%%%%%%% REFERÊNCIAS %%%%%%%%%%%%%%%%%

%\newpage
\phantomsection
\addcontentsline{toc}{section}{REFERÊNCIAS}

\nocite{*}
\bibliographystyle{ieeetr}
\bibliography{bib/bibliografia}

\newpage
\appendix

% Apêndice A
\include{apendice}

\end{document}
